The simulation of human attitudes through Large Language Models (LLMs) has
established a new paradigm in computational social science, termed
\textit{silicon sampling}~\cite{argyle_out_2023,filippas_large_2024}. Autoregressive
models encode latent socio-behavioral priors capable of reproducing collective
demographic viewpoints~\cite{argyle_out_2023,grossmann_ai_2023}, outperforming
classical classifiers in aggregate opinion alignment~\cite{miranda_simulating_2025}.
However, commercial foundation models disproportionately reflect educated, affluent,
and progressive Western strata~\cite{santurkar_whose_2023,bisbee_synthetic_2024},
underrepresenting Global South viewpoints~\cite{durmus_towards_2023} and
exhibiting persistent Anglo-American cultural skews in social and political value
judgments~\cite{hartmann_political_2023,masoud_cultural_2023,johnson_ghost_2022}.
Consequently, faithful simulation in non-WEIRD contexts requires culturally
grounded personas, structured prompt engineering, and empirical calibration.

Bridging demographic profiles with faithful decision-making demands disciplined
prompt design. While LLM execution can be conceptualized as persona role-play
over a generative distribution~\cite{shanahan_role_2023}, unconstrained persona
assignment risks caricaturization and stereotyping~\cite{deshpande_toxicity_2023,salewski_-context_2023}.
Grounding synthetic personas in official census microdata provides essential
empirical anchor points. To structure deliberative choices, Chain-of-Thought (CoT)
prompting elicits latent reasoning trajectories by producing intermediate tokens
prior to emitting final responses~\cite{wei_chain--thought_2022,kojima_large_2022},
reducing stochastic variance in multi-alternative opinion spaces~\cite{wang_self-consistency_2022}.
Furthermore, introducing contrasting calibration personas via few-shot examples
stabilizes response distributions against in-context recency and label
biases~\cite{brown_language_2020,zhao_calibrate_2021}.

Finally, to audit black-box agent behavior and verify persona consistency,
generative simulation incorporates principles of Explainable Artificial Intelligence
(XAI)~\cite{ribeiro_why_2016,guidotti_survey_2019}. Rather than treating models
as opaque predictors, qualitative inspection of intermediate reasoning
trajectories enables researchers to verify demographic semantic coherence,
detect persona cognitive dissonance, and explain divergent behavioral priors
between Brazilian Portuguese-centric and global baseline models.

